\section{Pilot deployment}\label{sec:evaluation}

\subsection{Setting and data}

The numbers below come from the production database of a single installation
used by the first author between April and June 2026. They are a
\emph{deployment record}, not a controlled experiment: the projects were
screened because they had to be screened, not because they were designed as an
evaluation. We report them because they establish the order of magnitude of
cost and scale, and because the journal's guidance asks for exactly this kind
of operational evidence alongside the agreement study
\autocite{Farotimi_2025}. The formal agreement study on a fresh blind sample is
described in Section~\ref{sec:planned} and is not yet complete.

Four projects are included. Three of them are screening runs over the same
source of 4{,}670 records under different prompt configurations, which we
label P1, P2 and P3; the fourth, P4, is an educational corpus of 343 records
used to exercise the platform end to end.

\subsection{Scale and decision distribution}

Table~\ref{tab:corpus} summarises the corpus. Of 14{,}353 records, 14{,}219
(99.1\%) carry a stored model decision. The 134 records without one are
concentrated in P1, the earliest run, and correspond to records imported after
the batch had been executed. The proportion of records classified as
\emph{doubt} falls sharply across the three runs over the same source, from
176 in P1 to 27 in P3, which is consistent with successive prompt revisions
that pushed the model towards a binary decision.

\begin{table}[htbp]
\centering
\caption{Records, model decisions and decision distribution per project}
\label{tab:corpus}
\begin{tabular}{@{}lrrrrrrr@{}}
\toprule
Project & Records & With AI & Included & Excluded & Doubt & Duplicates & Inclusion \\
 & & decision & & & & & rate \\
\midrule
P1 & 4{,}670 & 4{,}573 (97.9\%) & 565 & 3{,}899 & 176 & 30 & 12.1\% \\
P2 & 4{,}670 & 4{,}652 (99.6\%) & 913 & 3{,}660 & 79 & 18 & 19.6\% \\
P3 & 4{,}670 & 4{,}651 (99.6\%) & 1{,}030 & 3{,}594 & 27 & 19 & 22.1\% \\
P4 & 343 & 343 (100\%) & 11 & 332 & 0 & 0 & 3.2\% \\
\midrule
Total & 14{,}353 & 14{,}219 (99.1\%) & 2{,}519 & 11{,}485 & 282 & 67 & 17.6\% \\
\bottomrule
\end{tabular}
\end{table}

The inclusion rate is the number that matters for planning, and it varies by a
factor of almost two across configurations of the same protocol (12.1\% in P1
against 22.1\% in P3). A review team that estimates reviewer workload from a
single automated run would therefore be wrong by a wide margin. This is an
argument for treating the automated decisions as a first pass whose yield has
to be measured per project, which is what the doubt mechanism and the blind
mode are for.

\subsection{Cost per record}

Table~\ref{tab:cost} reports token consumption. The aggregate is 11.85 million
tokens for 14{,}219 records, or 833 tokens per record, with a range from 427 to
1{,}149 depending on prompt length and abstract length. The three tariff
scenarios translate this into money; they span two orders of magnitude because
provider prices differ by that much, and the platform records the model per
decision so that the applicable tariff can be reconstructed later.

\begin{table}[htbp]
\centering
\caption{Token consumption per project (top) and the resulting cost of the
whole corpus under three provider tariffs (bottom). Tariffs are expressed in US
dollars per million input and output tokens}
\label{tab:cost}
\small
\begin{tabular}{@{}lrrr@{}}
\toprule
Project & Input tokens & Output tokens & Tokens per record \\
\midrule
P1 & 3{,}846{,}046 & 279{,}519 & 902 \\
P2 & 5{,}032{,}476 & 313{,}751 & 1{,}149 \\
P3 & 1{,}656{,}524 & 333{,}570 & 427 \\
P4 & 367{,}003 & 21{,}739 & 1{,}133 \\
\midrule
Total & 10{,}902{,}049 & 948{,}579 & 833 \\
\bottomrule
\end{tabular}

\vspace{2mm}
\begin{tabular}{@{}lrr@{}}
\toprule
Tariff (input / output) & Cost of the corpus & US\$ per 1{,}000 records \\
\midrule
Low (0.10 / 0.40) & US\$1.47 & 0.10 \\
Mid (0.30 / 2.50) & US\$5.64 & 0.40 \\
High (1.25 / 10.00) & US\$23.11 & 1.63 \\
\bottomrule
\end{tabular}
\end{table}

The order of magnitude is the finding: screening a 15{,}000-record review costs
between one and twenty-five US dollars in model calls. The binding constraint
on adoption is therefore not the price of inference but the reliability of the
decisions, which is what the next subsection addresses.

\subsection{Agreement with the human reviewer}

Table~\ref{tab:kappa} reports every pair of records for which both an
automated and a stored human decision existed at the time of measurement. We
report all four projects, including the two whose agreement is close to
chance, because suppressing them would reproduce the reporting practice the
guidance warns against \autocite{Farotimi_2025}.

\begin{table}[htbp]
\centering
\caption{Agreement between automated and stored human decisions, with the
confusion matrix of the four projects (VN, true negative; FP, false positive;
VP, true positive; FN, false negative)}
\label{tab:kappa}
\small
\begin{tabular}{@{}lrrrrl@{}}
\toprule
Project & Pairs & $\kappa$ & 95\% CI & Raw agreement & Confusion matrix \\
\midrule
P1 & 4 & 0.000 & not computed ($n<5$) & 50.0\% & 2 VN, 2 FP, 0 VP, 0 FN \\
P2 & 34 & 0.084 & 0.000 to 0.235 & 47.1\% & 14 VN, 18 FP, 2 VP, 0 FN \\
P3 & 34 & 0.013 & $-0.290$ to 0.294 & 47.1\% & 6 VN, 14 FP, 10 VP, 4 FN \\
P4 & 31 & 0.608 & 0.290 to 0.903 & 83.9\% & 20 VN, 5 FP, 6 VP, 0 FN \\
\bottomrule
\end{tabular}
\end{table}

Three observations follow, and the third is the one we consider most important.

First, only P4 supports an interpretation. With 31 pairs, $\kappa$ is 0.61
(95\% CI 0.29 to 0.90), raw agreement 83.9\%, sensitivity 100\% (six of six
included records recovered) and specificity 80\% (20 of 25 exclusions
confirmed). The interval is wide enough that the point estimate should not be
quoted as the accuracy of the tool; it is a signal, not a measurement.

Second, the error pattern in P2 and P3 is asymmetric and consistent across both
runs: over-inclusion, with 14 to 18 false positives against 0 to 4 false
negatives. In screening, a false negative costs evidence and a false positive
costs reviewer time. The configuration used in these runs is therefore
recall-oriented, which is defensible for a first pass but has a direct
consequence for workload: the reviewer still reads a substantial share of the
records, and the workload reduction is smaller than the raw automation rate
suggests.

Third, the low agreement in P2 and P3 is not evidence that the model cannot
screen. The stored human decisions in those projects were recorded during
development, at dates that were not captured, and may have been made under an
earlier version of the criteria. What the table demonstrates is a limitation of
the measurement, not of the model, and it is precisely the limitation that the
blind validation mode was built to remove: the platform now records who
decided, when, and against which protocol version, which the pilot data do not
carry for every pair.

\subsection{Engineering quality}

The platform's own test suite is reported because reproducibility of a
research tool is part of its evaluation. At the time of the pilot the backend
suite contained 260 tests, of which 259 passed and one was skipped, in 2
minutes and 42 seconds, with 75.11\% line coverage. Throughput during batch
execution produced a median interval between consecutive decision events of
0.00 to 0.05 seconds per project, confirming that screening runs in a single
batch rather than interactively; it is not a measure of reviewer time.

\subsection{What has not been measured}\label{sec:planned}

The guidance for generative AI submissions requires workload metrics for tools
whose purpose is to reduce effort \autocite{Farotimi_2025}. Sifter has not yet
produced them, and we state the gap rather than approximate it. Three
quantities are outstanding.

\begin{enumerate}
\item \textbf{Work saved over screening (WSS)} and \textbf{number needed to
read (NNR)}, computed on a blind sample of 300 pairs drawn from the
educational corpus. The sample and the reference decisions have been generated
and are stored; the reference screening has not been completed.
\item \textbf{Reviewer time per record}, with and without the automated first
pass, which requires a timed screening session and is the input for any
statement about hours saved.
\item \textbf{External validation} on a review protocol and record set produced
by another team, which is the only way to test transfer across fields and
prompt styles.
\end{enumerate}

Until these three exist, the claim this paper supports is bounded: the
platform records what it does, at a cost that is negligible relative to
reviewer time, and it makes the agreement study cheap enough to be run as
routine quality control. It does not yet support a claim about accuracy in
general.
